\subsection{紧算子的不变子空间}

应当把下面的定理与舒尔引理（A, VIII, p. 43, cor. 及 V, p. 386, prop. 6）以及 I, p. 26 的推论 4 相比较。

定理 3. --- 设 E 是 C 上的分离局部凸空间，A 是 \(\mathcal{L}(\mathrm{E})\) 的一个子集。我们作下列假设：

\begin{enumerate}
\def\labelenumi{(\roman{enumi})}
\item
  不存在 E 的既异于 \(\{0\}\) 又异于 E 的、在 A 下稳定的闭向量子空间；
\item
  集合 A 含有一个非零的紧自同态。那么 A 的交换子化归于恒等映射的纯量倍数。
\end{enumerate}

用 A 生成的 \(\mathcal{L}(\mathrm{E})\) 的含单位元子代数代替 A，就化归为 A 是 \(\mathcal{L}(\mathrm{E})\) 的一个含单位元子代数的情形。于是设 \(h\) 是属于 A 的一个非零紧自同态。

引理 1. --- 存在 A 的一个元素 a，使得 \(ha - 1_{E}\) 的核不化归于 0。

存在 E 的一个元素 \(x_{0}\)，使得 \(h(x_{0}) \neq 0\)。设 V 是 \(h(x_{0})\) 的一个不包含 0 的闭邻域。由于自同态 h 紧，可以取 \(x_{0}\) 的一个开凸邻域 U，使得 \(h(\mathrm{U})\) 是 V 的一个相对紧子集。于是 \(h(\mathrm{U})\) 的闭包 C 是 E 的一个紧凸子集；由于它包含于 V，它不包含 0。

设 \(y\) 是 C 的一个点。用 \(Ay\) 表示 \(y\) 在 A 的诸元素下的像所成的集。它是 E 的一个在 A 下稳定的向量子空间；由于它包含 \(y\)，它是非零的。由定理的假设 (i)，集合 \(Ay\) 在 E 中稠密。因此存在 A 的一个元素 \(b\)，使得 \(b(y) \in U\)。由于集合 C 紧，存在 A 中元素的一个有限族 \((b_1, \ldots, b_n)\)，使得诸集合 \(b_i^{-1}(U)\) 覆盖 C。存在 C 上从属于这个覆盖的一个连续单位分解 \((\varphi_1, \ldots, \varphi_n)\)（TG, IX, p. 47, th. 3）。我们用下式定义映射 \(f: C \to E\)：

\[
f (x) = \sum_ {i = 1} ^ {n} \varphi_ {i} (x) b _ {i} (x)
\]

其中 \(x \in \mathrm{C}\)；这是一个连续映射。由于 U 是凸的，且 \(b_i(x)\) 在 \(\varphi_i(x)\) 非零时属于 U，我们有 \(f(\mathrm{C}) \subset \mathrm{U}\) 及 \(h(f(\mathrm{C})) \subset \mathrm{C}\)。由勒雷--绍德尔定理（III, p. 11, 定理 2），存在一个元素 \(x \in \mathrm{C}\)，使得 \(h(f(x)) = x\)。于是令

\[
a = \sum_ {i = 1} ^ {n} \varphi_ {i} (x) b _ {i}.
\]

我们有 \(h(a(x)) = h(f(x)) = x\)，且由于 0 不属于 C，x 非零，因此 \(ha - 1_{E}\) 的核非零。

我们来完成定理 3 的证明。设 \(a\) 是 A 的一个元素，使得 \(ha - 1_{\mathrm{E}}\) 的核 T 非零（引理 1）。用 \(i\) 表示 T 到 E 中的典范单射。线性映射 \(i\) 等于 \(h \circ a \circ i\)，因此它是紧的。于是 T 的维数有限（III, p. 2, 注记 3）。设 \(u\) 是 \(\mathcal{L}(\mathrm{E})\) 的一个与 A 交换的元素。由于 \(u\) 与 \(h \circ a\) 交换，我们有 \(u(\mathrm{T}) \subset \mathrm{T}\)。由于 T 在代数闭域 \(\mathbf{C}\) 上是有限维且非零的，\(u\) 在 T 上诱导的自同态有一个特征值 \(\lambda\)。

令 \(\mathrm{F} = \mathrm{Ker}(u - \lambda1_{\mathrm{E}})\)。这是 E 的一个非零闭向量子空间；由于 u 与 A 的诸元素交换，它在 A 下稳定。由假设 (i)，我们有 F = E，从而 \(u = \lambda1_{E}\)。

推论 1. --- 保持定理 3 的假设，并再设 A 的诸元素两两交换。那么 E 在 C 上是一维的。

由定理 3 的假设 (ii)，空间 E 非零。因此它包含一个维数为 1 的向量子空间 F。由于 E 假设为分离的，这个子空间是闭的。由定理 3，A 的每个元素都是位似变换，故使 F 稳定。于是由定理 3 的假设 (i) 推出 F 等于 E。

推论 2（洛莫诺索夫）. --- 设 E 是域 C 上维数至少为 2 的分离局部凸空间，u 是 E 的一个自同态。假设存在 E 的一个与 u 交换的紧自同态 \(h \neq 0\)。那么存在 E 的一个既异于 \(\{0\}\) 又异于 E 的闭向量子空间 F，使得 \(u(F) \subset F\)。

事实上，推论 1 表明集合 A = \{u, h\} 不可能满足定理 3 的假设 (i)。

推论 3. --- 设 E 是域 C 上的分离局部凸空间，u 是 E 的紧自同态。若 E 的维数至少为 2，则存在 E 的一个既异于 \(\{0\}\) 又异于 E 的闭向量子空间 F，使得 \(u(F) \subset F\)。

u = 0 的情形是直接的。\(u \neq 0\) 的情形由推论 2 取 h = u 得出。

